\section{A further sixth-root improvement}\label{app:sixth-improvement}
The recovered manuscript gave $Q\ge9$. The following additional finite exclusion improves that bound.
\begin{theorem}\label{thm:sixth-ten}
Every $H\in\mu_6^{15\times15}$ satisfies $Q(H)\ge10$. Consequently
\[5\cdot2^{26}\sqrt3\le D_6(15)\le
 \left[(15-\sqrt{2/21})^{14}(15+14\sqrt{2/21})\right]^{1/2}.\]
\end{theorem}
\begin{proof}
The alphabet-independent frame gap and the odd norm obstruction leave only $Q=9$ to exclude. Put $T=HH^*-15I$, $S=H^*H-15I$. Their entries are Eisenstein integers and every nonzero edge costs at least one. An isolate on one side forbids leaves on the other, since an isolated row restricted to an opposite component is a flat unit kernel vector. Hence a support cannot be isolate-free while the other has isolates: it would be leafless on fifteen vertices and require at least fifteen edges.
If both supports have isolates, their cores are leafless and have at most nine vertices. Enumerate every connected leafless support with at most nine edges: the graph atlas covers orders at most seven; orders eight and nine are generated from every nonisomorphic spanning tree and every admissible added edge set. There are $100$ supports. Enumerate every Eisenstein norm allocation of total at most nine, fix tree phases by unit switching, and assign every chord phase. This gives $51{,}774$ labelled states. A flat kernel vector requires a positive diagonal in the kernel projector. Grouping equal spectra and retaining the largest observed minimum projector diagonal is a relaxation, leaving $365$ component records and $285$ energy-nine assemblies.
For a component $C$ and $r'$ opposite isolates, the kernel projector must satisfy $(P_{\ker T_C})_{ii}\ge r'/15$ for every $i$: the projection onto those orthogonal columns, divided by fifteen, has this constant diagonal and lies in the kernel. After the norm sieve and equal-spectrum pairing, twenty-nine pairs remain. Each violates one of these rational inequalities. Every projector is computed as $q(T_C)/q(0)$, where $q$ is the characteristic polynomial with its zero factors removed. All comparisons are exact.
If neither support has isolates, nine energy units force nine edges, all of norm one. The support is either a forest with component-size partition
\[(5,2,2,2,2,2),\ (4,3,2,2,2,2),\ (3,3,3,2,2,2),\]
or a triangle and six pairs. Tree phases switch away. Every forest determinant has odd five-valuation, verified directly from the finitely many trees of orders two through five. For the triangle, write its flux trace as $\tau\in\{-2,-1,1,2\}$. Its total determinant is $(3330+\tau)224^6$. The norm test excludes $\tau=-1,2$. The remaining triangle polynomials are $x^3-3x+2=(x-1)^2(x+2)$ and $x^3-3x-1$. Against the twelve pair coordinates, the spectral rectangle areas are $36$, while the capacities are respectively $32$ and $0$. Both are impossible. Substituting $Q=10$ into the envelope proves the upper bound.
\end{proof}
The certificate is regenerated by \file{extensions/sixth_roots/verify_gap.py}. It uses neither a determinant threshold nor the third-root orthogonality capacity. The exhibited sixth-root matrix still has energy $53$; the new bound is an obstruction, not a maximum-determinant classification.
